\documentclass{beamer}

\usetheme{Madrid}
\definecolor{unacblue}{RGB}{0,51,102}
\setbeamercolor{structure}{fg=unacblue}

\usepackage{graphicx}
\usepackage{tikz}
\usepackage[spanish]{babel}
\usepackage{amsmath}

\setbeamertemplate{navigation symbols}{}

\title{Estadística Descriptiva:\\Scatter plot, covarianza y coeficiente de correlación}
\author{Est. Luis Ayrton Vega Ayquipa \\
\small lavegaa@unac.edu.pe}
\institute{Universidad Nacional del Callao \\
Facultad de Economía \\
Escuela Profesional de Economía}
\date{\today}

% Pie de página personalizado
\setbeamertemplate{footline}{
  \leavevmode%
  \hbox{%
    \begin{beamercolorbox}[wd=.33\paperwidth,ht=2.5ex,dp=1ex,left]{author in head/foot}
      \hspace{1em}UNAC FCE - LAVA  
    \end{beamercolorbox}%
    \begin{beamercolorbox}[wd=.34\paperwidth,ht=2.5ex,dp=1ex,center]{title in head/foot}
      Ayudantía de Estadística Descriptiva 2026 A 
    \end{beamercolorbox}%
    \begin{beamercolorbox}[wd=.33\paperwidth,ht=2.5ex,dp=1ex,right]{date in head/foot}
      \insertshortdate{} \hspace{1em} \insertframenumber/\inserttotalframenumber\hspace{1em}
    \end{beamercolorbox}%
  }%
  \vskip0pt%
}


\begin{document}

%====================================================
\begin{frame}
\titlepage
\end{frame}

%====================================================
\begin{frame}{¿Qué es un Scatter Plot?}

\textbf{Definición:}

El diagrama de dispersión es una representación gráfica de pares ordenados:

\[
(X,Y)
\]

\vspace{0.3cm}

\begin{itemize}
\item Cada punto representa una observación.
\item Permite visualizar relaciones entre variables.
\item Es una herramienta fundamental del:
\[
EDA \quad (\text{Exploratory Data Analysis})
\]
\end{itemize}

\end{frame}

%====================================================
\begin{frame}{¿Para qué sirve?}

\begin{itemize}
\item Detectar relaciones lineales
\item Identificar patrones
\item Detectar valores atípicos
\item Analizar dirección de asociación
\item Evaluar dispersión de datos
\end{itemize}

\vspace{0.4cm}

\begin{block}{Importancia en econometría}
Antes de estimar una regresión, es necesario observar cómo se relacionan las variables.
\end{block}

\end{frame}

%====================================================
\begin{frame}{Interpretación básica}

\begin{columns}

\column{0.5\textwidth}

\textbf{Relación positiva}

\begin{itemize}
\item Cuando \(X\) aumenta
\item \(Y\) también aumenta
\end{itemize}

\vspace{0.4cm}

\textbf{Ejemplo:}
\begin{itemize}
\item educación e ingresos
\end{itemize}

\column{0.5\textwidth}

\textbf{Relación negativa}

\begin{itemize}
\item Cuando \(X\) aumenta
\item \(Y\) disminuye
\end{itemize}

\vspace{0.4cm}

\textbf{Ejemplo:}
\begin{itemize}
\item precio y demanda
\end{itemize}

\end{columns}

\end{frame}

%====================================================
\begin{frame}{Ejemplo aplicado}

\textbf{Pregunta económica:}

\vspace{0.2cm}

¿Cómo influye el número de horas de estudio sobre la nota obtenida?

\vspace{0.4cm}

\begin{center}
\begin{tabular}{cc}
\toprule
Horas de estudio (\(X\)) & Nota (\(Y\)) \\
\midrule
2 & 10 \\
3 & 11 \\
4 & 13 \\
5 & 15 \\
6 & 16 \\
7 & 18 \\
\bottomrule
\end{tabular}
\end{center}

\end{frame}

%====================================================
\begin{frame}{Construcción del Scatter Plot}

\textbf{Paso 1:}

Definir variables:

\[
X = \text{Horas de estudio}
\]

\[
Y = \text{Nota}
\]

\vspace{0.4cm}

\textbf{Paso 2:}

Representar cada observación como un punto:

\[
(2,10), (3,11), (4,13)
\]

etc.

\end{frame}

%====================================================
\begin{frame}{Scatter Plot del ejemplo}

\begin{center}

\begin{tikzpicture}

\begin{axis}[
width=10cm,
height=6cm,
xlabel={Horas de estudio},
ylabel={Nota},
grid=major
]

\addplot[
only marks,
mark=*,
]
coordinates {
(2,10)
(3,11)
(4,13)
(5,15)
(6,16)
(7,18)
};

\end{axis}

\end{tikzpicture}

\end{center}

\end{frame}

%====================================================
\begin{frame}{Interpretación del gráfico}

\begin{itemize}
\item Existe una relación positiva.
\item A mayor número de horas de estudio:
\[
\Rightarrow
\]
mayor nota promedio.
\item Los puntos siguen aproximadamente una tendencia lineal.
\end{itemize}

\vspace{0.4cm}

\begin{block}{Conclusión}
El scatter plot sugiere que podría existir correlación positiva entre ambas variables.
\end{block}

\end{frame}

%====================================================
\begin{frame}{Scatter Plot y EDA}

\textbf{En análisis exploratorio de datos (EDA), el scatter plot permite:}

\vspace{0.3cm}

\begin{itemize}
\item Explorar relaciones preliminares
\item Verificar linealidad
\item Detectar heterogeneidad
\item Identificar observaciones atípicas
\item Decidir si aplicar regresión lineal
\end{itemize}

\vspace{0.4cm}

\begin{block}{¡Importante!}
El scatter plot es el primer paso visual antes del modelamiento econométrico.
\end{block}

\end{frame}

%====================================================
\begin{frame}{Conclusión}

\begin{itemize}
\item El scatter plot representa pares ordenados.
\item Permite analizar relaciones entre variables.
\item Es fundamental en estadística y econometría.
\item Ayuda a comprender la estructura de los datos antes de estimar modelos.
\end{itemize}

\end{frame}

\begin{frame}{¿Qué es la covarianza?}

\textbf{Definición:}

La covarianza mide cómo dos variables cambian conjuntamente.

\vspace{0.4cm}

\begin{itemize}
\item Si ambas variables aumentan juntas:
\[
Cov(X,Y)>0
\]

\item Si una aumenta y la otra disminuye:
\[
Cov(X,Y)<0
\]

\item Si no existe relación lineal:
\[
Cov(X,Y)\approx0
\]
\end{itemize}

\end{frame}

%====================================================
\begin{frame}{Interpretación económica}

\begin{block}{Ejemplos}
\begin{itemize}
\item Educación e ingresos:
\[
Cov(X,Y)>0
\]

\item Precio y demanda:
\[
Cov(X,Y)<0
\]
\end{itemize}
\end{block}

\vspace{0.3cm}

\textbf{Idea principal:}

La covarianza permite identificar la dirección de la relación entre variables.

\end{frame}

%====================================================
\begin{frame}{Suma producto de desviaciones}

\textbf{Concepto fundamental:}

La covarianza se basa en:

\[
(x_i-\bar{x})(y_i-\bar{y})
\]

\vspace{0.1cm}

\textbf{Interpretación:}

\begin{itemize}
\item Mide cómo se desvían conjuntamente las observaciones.
\item Compara cada dato respecto a sus medias.
\end{itemize}

\vspace{0.0cm}

\begin{block}{Casos}
\begin{itemize}
\item Ambos positivos:
\[
(+)(+) = +
\]

\item Ambos negativos:
\[
(-)(-) = +
\]

\item Signos distintos:
\[
(+)(-) = -
\]
\end{itemize}
\end{block}

\end{frame}

%====================================================
\begin{frame}{Fórmula de covarianza poblacional}

\[
\sigma_{XY}
=
\frac{\sum (x_i-\mu_X)(y_i-\mu_Y)}{N}
\]

\vspace{0.4cm}

\textbf{Donde:}

\begin{itemize}
\item \(\mu_X\): media poblacional de \(X\)
\item \(\mu_Y\): media poblacional de \(Y\)
\item \(N\): número de observaciones
\end{itemize}

\end{frame}

%====================================================
\begin{frame}{Fórmula de covarianza muestral}

\[
S_{XY}
=
\frac{\sum (x_i-\bar{x})(y_i-\bar{y})}{n-1}
\]

\vspace{0.4cm}

\textbf{Diferencia importante:}

\begin{itemize}
\item Población:
\[
N
\]

\item Muestra:
\[
n-1
\]
\end{itemize}

\vspace{0.3cm}

\textbf{¿Por qué \(n-1\)?}

Porque se pierde un grado de libertad al estimar la media.

\end{frame}

%====================================================
\begin{frame}{Ejemplo aplicado}

\textbf{Relación entre horas de estudio y nota}

\vspace{0.3cm}

\begin{center}
\begin{tabular}{cc}
\toprule
Horas (\(X\)) & Nota (\(Y\)) \\
\midrule
2 & 10 \\
3 & 11 \\
4 & 13 \\
5 & 15 \\
6 & 16 \\
\bottomrule
\end{tabular}
\end{center}

\vspace{0.3cm}

\[
\bar{x}=4
\]

\[
\bar{y}=13
\]

\end{frame}

%====================================================
\begin{frame}{Tabla de cálculo}

\small

\begin{center}
\begin{tabular}{cccccc}
\toprule
\(X_i\) & \(Y_i\) &
\(X_i-\bar{x}\) &
\(Y_i-\bar{y}\) &
Producto \\
\midrule

2 & 10 & -2 & -3 & 6 \\
3 & 11 & -1 & -2 & 2 \\
4 & 13 & 0 & 0 & 0 \\
5 & 15 & 1 & 2 & 2 \\
6 & 16 & 2 & 3 & 6 \\

\bottomrule
\end{tabular}
\end{center}

\vspace{0.3cm}

\[
\sum (x_i-\bar{x})(y_i-\bar{y}) = 16
\]

\end{frame}

%====================================================
\begin{frame}{Cálculo de la covarianza}

Usando la fórmula muestral:

\[
S_{XY}
=
\frac{\sum (x_i-\bar{x})(y_i-\bar{y})}{n-1}
\]

\vspace{0.4cm}

Sustituyendo:

\[
S_{XY}
=
\frac{16}{5-1}
\]

\[
S_{XY}=4
\]

\end{frame}

%====================================================
\begin{frame}{Interpretación del resultado}

\[
S_{XY}=4
\]

\vspace{0.4cm}

\textbf{Interpretación:}

\begin{itemize}
\item La covarianza es positiva.
\item Existe relación directa entre:
\begin{itemize}
\item horas de estudio
\item nota obtenida
\end{itemize}

\item Cuando aumentan las horas:
\[
\Rightarrow
\]
las notas tienden a aumentar.
\end{itemize}

\end{frame}

%====================================================
\begin{frame}{Limitación de la covarianza}

\begin{block}{Problema}
La covarianza depende de las unidades de medida.
\end{block}

\vspace{0.4cm}

\textbf{Ejemplo:}

\begin{itemize}
\item Horas y notas
\item Soles y consumo
\item Edad y salario
\end{itemize}

\vspace{0.3cm}

Por ello se utiliza posteriormente:

\[
\text{Coeficiente de correlación}
\]

\end{frame}

%====================================================
\begin{frame}{Conclusión}

\begin{itemize}
\item La covarianza mide variación conjunta.
\item Se basa en productos de desviaciones.
\item Permite identificar dirección de relación.
\item Es fundamental para correlación y regresión.
\end{itemize}

\end{frame}

%====================================================
\begin{frame}{¿Qué es el coeficiente de correlación?}

\textbf{Definición:}

El coeficiente de correlación mide:

\begin{itemize}
\item dirección de la relación
\item intensidad de la relación lineal
\end{itemize}

entre dos variables.

\vspace{0.4cm}

\[
-1 \leq r \leq 1
\]

\vspace{0.3cm}

\begin{block}{Interpretación}
Mientras más cercano sea el valor a \(\pm1\), más fuerte será la relación lineal.
\end{block}

\end{frame}

%====================================================
\begin{frame}{Interpretación del coeficiente}

\begin{center}

\begin{tabular}{cc}
\toprule
Valor de \(r\) & Interpretación \\
\midrule
\(r=1\) & Relación positiva perfecta \\
\(0<r<1\) & Relación positiva \\
\(r=0\) & Sin relación lineal \\
\(-1<r<0\) & Relación negativa \\
\(r=-1\) & Relación negativa perfecta \\
\bottomrule
\end{tabular}

\end{center}

\vspace{0.4cm}

\textbf{Ejemplos económicos:}

\begin{itemize}
\item Educación e ingresos:
\[
r>0
\]

\item Precio y demanda:
\[
r<0
\]
\end{itemize}

\end{frame}

%====================================================
\begin{frame}{¿Por qué surge la correlación?}

La covarianza tiene un problema:

\vspace{0.3cm}

\begin{block}{Problema}
La covarianza depende de las unidades de medida.
\end{block}

\vspace{0.4cm}

Por ello se estandariza utilizando desviaciones estándar.

\vspace{0.4cm}

Así obtenemos una medida:

\begin{itemize}
\item sin unidades
\item comparable
\item acotada entre -1 y 1
\end{itemize}

\end{frame}

%====================================================
\begin{frame}{Fórmula poblacional}

\[
\rho_{XY}
=
\frac{
Cov(X,Y)
}{
\sigma_X \sigma_Y
}
\]

\vspace{0.4cm}

\textbf{Donde:}

\begin{itemize}
\item \(\rho_{XY}\): correlación poblacional
\item \(Cov(X,Y)\): covarianza poblacional
\item \(\sigma_X\): desviación estándar de \(X\)
\item \(\sigma_Y\): desviación estándar de \(Y\)
\end{itemize}

\end{frame}

%====================================================
\begin{frame}{Fórmula muestral}

\[
r
=
\frac{
S_{XY}
}{
S_X S_Y
}
\]

\vspace{0.4cm}

\textbf{Forma desarrollada:}

\[
r=
\frac{
\sum (x_i-\bar{x})(y_i-\bar{y})
}{
\sqrt{
\sum (x_i-\bar{x})^2
\sum (y_i-\bar{y})^2
}
}
\]

\end{frame}

%====================================================
\begin{frame}{¿Qué es una suma de cuadrados?}

\textbf{Concepto:}

Una suma de cuadrados mide variabilidad.

\vspace{0.4cm}

\[
\sum (x_i-\bar{x})^2
\]

\vspace{0.4cm}

\textbf{Interpretación:}

\begin{itemize}
\item Mide qué tan dispersos están los datos respecto a la media.
\item Es la base de:
\begin{itemize}
\item varianza
\item desviación estándar
\item correlación
\item regresión
\end{itemize}
\end{itemize}

\end{frame}

%====================================================
\begin{frame}{Ejemplo muestral}

\textbf{Relación entre horas de estudio y nota}

\vspace{0.3cm}

\begin{center}
\begin{tabular}{cc}
\toprule
Horas (\(X\)) & Nota (\(Y\)) \\
\midrule
2 & 10 \\
3 & 11 \\
4 & 13 \\
5 & 15 \\
6 & 16 \\
\bottomrule
\end{tabular}
\end{center}

\vspace{0.3cm}

\[
\bar{x}=4
\]

\[
\bar{y}=13
\]

\end{frame}

%====================================================
\begin{frame}{Tabla de cálculo}

\small

\begin{center}
\begin{tabular}{ccccccc}
\toprule

\(X_i\) &
\(Y_i\) &
\(X_i-\bar{x}\) &
\(Y_i-\bar{y}\) &
\((X_i-\bar{x})^2\) &
\((Y_i-\bar{y})^2\) &
Producto \\

\midrule

2 & 10 & -2 & -3 & 4 & 9 & 6 \\
3 & 11 & -1 & -2 & 1 & 4 & 2 \\
4 & 13 & 0 & 0 & 0 & 0 & 0 \\
5 & 15 & 1 & 2 & 1 & 4 & 2 \\
6 & 16 & 2 & 3 & 4 & 9 & 6 \\

\bottomrule
\end{tabular}
\end{center}

\vspace{0.3cm}

\[
\sum (x_i-\bar{x})^2 = 10
\]

\[
\sum (y_i-\bar{y})^2 = 26
\]

\[
\sum (x_i-\bar{x})(y_i-\bar{y}) = 16
\]

\end{frame}

%====================================================
\begin{frame}{Cálculo del coeficiente de correlación}

Aplicando:

\[
r=
\frac{
\sum (x_i-\bar{x})(y_i-\bar{y})
}{
\sqrt{
\sum (x_i-\bar{x})^2
\sum (y_i-\bar{y})^2
}
}
\]

\vspace{0.4cm}

Sustituyendo:

\[
r=
\frac{16}{\sqrt{(10)(26)}}
\]

\[
r=
\frac{16}{16.1245}
\]

\[
r \approx 0.992
\]

\end{frame}

%====================================================
\begin{frame}{Interpretación del resultado}

\[
r \approx 0.992
\]

\vspace{0.4cm}

\textbf{Interpretación:}

\begin{itemize}
\item Existe una correlación positiva muy fuerte.
\item A mayor número de horas de estudio:
\[
\Rightarrow
\]
mayor nota obtenida.
\item Los datos presentan comportamiento casi lineal.
\end{itemize}

\end{frame}

%====================================================
\begin{frame}{Ejemplo poblacional}

Supongamos una población de trabajadores:

\vspace{0.3cm}

\begin{center}
\begin{tabular}{cc}
\toprule
Experiencia (\(X\)) & Salario (\(Y\)) \\
\midrule
1 & 1200 \\
2 & 1400 \\
3 & 1700 \\
4 & 1900 \\
\bottomrule
\end{tabular}
\end{center}

\vspace{0.4cm}

Aplicando:

\[
\rho_{XY}
=
\frac{
Cov(X,Y)
}{
\sigma_X \sigma_Y
}
\]

\vspace{0.3cm}

se obtiene:

\[
\rho_{XY}\approx0.98
\]

\end{frame}

%====================================================
\begin{frame}{Importancia del coeficiente de correlación}

\begin{itemize}
\item Detectar relaciones lineales
\item Analizar intensidad de asociación
\item Base de regresión lineal
\item Herramienta central en econometría
\end{itemize}

\vspace{0.4cm}

\begin{block}{Importante}
Correlación no implica causalidad.
\end{block}

\end{frame}

%====================================================
\begin{frame}{Conclusión}

\begin{itemize}
\item La correlación mide dirección e intensidad.
\item Surge de estandarizar la covarianza.
\item Utiliza sumas de cuadrados.
\item Es fundamental para regresión y análisis económico.
\end{itemize}

\end{frame}



\end{document}